\chapter{NETRA on site}
\label{ch:i-netra}

At the site the Mac joins the robot's Wi-Fi and runs NETRA without internet; the robot connects
to it directly.

\section{Start NETRA}

\begin{lstlisting}[language=bash]
cd ~/x30/Netra
scripts/local.sh start     # the images built before: no internet needed
scripts/local.sh awake     # keep the Mac awake for schedules and the robot
scripts/local.sh lan       # the Mac's address on each network, e.g. http://192.168.1.50
\end{lstlisting}

Reserve that address for the Mac in the robot's router, or check it after each join: the robot
uses it.

\section{Connect the robot}

On the navigation host, check that NETRA answers, then give the robot its address, its ID and
its token:

\begin{lstlisting}[language=bash]
curl http://192.168.1.50/healthz                     # {"status":"ok"}
sudo nano /etc/x30/navigation.env
\end{lstlisting}

\begin{lstlisting}[language=bash]
X30_LAUNCH_ARGS="site:=gardu netra_url:=http://192.168.1.50 netra_robot_id:=x30-01"
X30_NETRA_TOKEN=<CC_TOKEN>
\end{lstlisting}

\begin{lstlisting}[language=bash]
sudo systemctl restart x30-navigation
journalctl -u x30-navigation -f | grep -i netra      # "command center": connected
\end{lstlisting}

Within seconds the robot is online on the dashboard: battery, mode, the navigation state, the
recorder and the inspection in \menupath{Live control}. \code{site:=} must be there: NETRA's
routes and their inspection plans go into the active site (\file{~/x30/map/<site>/netra/}).

\section{What the robot sends and keeps}

\begin{itemize}
  \item \textbf{Telemetry} every second; patrol and stop events, alerts and the answers to
    commands, all over one WebSocket that comes back by itself.
  \item \textbf{Uploads}: every taught route and every inspection of NETRA's points wait in a
    spool on the robot (\file{~/x30/spool}) until NETRA has them, across restarts and while NETRA
    is away.
  \item \textbf{Clock}: the robot has no internet clock; NETRA measures the offset and puts every
    measurement on its own clock.
\end{itemize}

\begin{warnbox}[Install the license first]
Without the substation module NETRA accepts an inspection upload but drops it, and the robot
does not send it again. The reports stay on the robot (\file{~/x30/reports}), but NETRA will
not have them.
\end{warnbox}
